\SourceLeaf{229}%
\bigskip\Needspace{9\baselineskip}
\pdfbookmark[1]{Ceres Ferdinandea. June 1802}{ceres180206}
\begin{center}
CERES FERDINANDEA.\par
\medskip\hrule\medskip
Monatliche Correspondenz. Vol. V. pp.\ 576—590. June 1802.\par
\medskip\hrule
\end{center}\medskip
—\quad —\quad —

Dr.\ \textsc{Gauss} too has continued, as before, to compare all our observations of this planet with elements VII of its orbit. This comparison confirms that the errors in these elements change very little and do so slowly, so the elements give the planet’s position with more accuracy than is needed throughout this year’s period of visibility. The comparison will also be of great use to him later, when he refines the elements further and takes the perturbation equations into account. Up to 19 March Dr.\ \textsc{Gauss} had carried out this comparison in the May issue, p.\ 470 [p.\ 211 of this volume]. Here are the results for the next twelve observations.
\begin{center}
{\small\setlength{\tabcolsep}{4pt}
\begin{tabular}{r r|r|r|r|r|r}
\multicolumn{2}{c|}{1802}&\shortstack{Mean time\\at Seeberg}&\shortstack{Calculated R.A.\\of Ceres}&\shortstack{Calculated\\declination\\of Ceres}&\multicolumn{2}{c}{\shortstack{Difference\\R—B}}\\
\multicolumn{2}{c|}{}&&&&in R.A.&in decl.\\\hline
March&23&$12^{\mathrm{u}}\quad6^{\mathrm{m}}\ 28^{\mathrm{s}}9$&$182^0\ 20'\ 39''9$&$17^0\ 36'\ 51''4$&$+16''1$&$+23''5$\\
&27&$11\ 47\ 20.3$&$181\ 29\ 21.2$&$17\ 50\ 55.5$&$+20.1$&$+25.9$\\
&28&$11\ 42\ 33.7$&$181\ 16\ 41.4$&$17\ 53\ 52.9$&$+23.7$&$+29.8$\\
&29&$11\ 37\ 48.0$&$181\quad4\quad7.5$&$17\ 56\ 36.3$&$+18.3$&$+32.3$\\
&30&$11\ 33\quad2.5$&$180\ 51\ 40.3$&$17\ 59\quad5.5$&$+16.4$&$+36.5$\\
&31&$11\ 28\ 17.0$&$180\ 39\ 21.1$&$18\quad1\ 20.3$&$+24.1$&$+31.9$\\
April&1&$11\ 23\ 32.4$&$180\ 27\ 10.7$&$18\quad3\ 20.6$&$+25.2$&$+30.8$\\
&2&$11\ 18\ 48.7$&$180\ 15\ 10.1$&$18\quad5\quad6.2$&$+23.6$&$+31.4$\\
&3&$11\ 14\quad6.0$&$180\quad3\ 20.1$&$18\quad6\ 37.4$&$+17.7$&$+27.0$\\
&4&$11\quad9\ 24.1$&$179\ 51\ 41.7$&$18\quad7\ 53.5$&$+10.8$&$+24.1$\\
&5&$11\quad4\ 41.9$&$179\ 40\ 15.6$&$18\quad8\ 54.9$&$+21.5$&$+27.6$\\
&7&$10\ 55\ 21.3$&$179\ 18\quad3.9$&$18\ 10\ 12.5$&$+24.2$&$+25.4$
\end{tabular}}
\end{center}

Dr.\ \textsc{Gauss} calculated the planet’s opposition to the Sun from our observations of 15, 16, 17, 18 and 19 March. He found that it occurred on 17 March at $4^{\mathrm{u}}\ 21^{\mathrm{m}}\ 10^{\mathrm{s}}$ mean Seeberg time, at longitude $5^{\mathrm{s}}\ 26^0\ 21'\ 27''3$ and geocentric latitude $17^0\quad7'\ 59''6$.
\par\noindent —\quad —\quad —\hfill [\textsc{von Zach}.]\par

